\chapter{Known Limitations}
\label{ch:limitations}

Stated in one place, as an honest list, not a formality --- each item below is
individually documented at the point in this report where it arises; this chapter
exists so a reader never has to reconstruct the full list by re-reading every chapter.

\begin{enumerate}
\item \textbf{Reduced-eye geometry throughout} (Section~\ref{ch:optical_system}): no
  anterior-chamber distance, no cornea/lens power split, single lumped refracting
  plane. Standard, but every pupil/lens-plane distance in this project's pipeline is a
  modeling convenience, not an anatomical one.
\item \textbf{No chromatic aberration, higher-order aberrations, Stiles--Crawford
  apodization in the core PSF, or finite aperture-membrane thickness}
  (Section~\ref{sec:unmodeled}).
\item \textbf{$r_{EE}(P)$ normalizes to grid-captured power, not emitted power}, and is
  measured over a different spatial domain than the threshold-crossing radii
  (Section~\ref{sec:effective-radius}). Both are quantified, not merely flagged, at
  this project's tested configurations.
\item \textbf{$r_{EE}(50\%)$, and at low myopia even $r_{EE}(70\%)$/$r_{EE}(80\%)$, can
  saturate at the simulation grid's own resolution floor}
  (Chapter~\ref{ch:percentile}) --- identified and quantified in this report, not
  previously documented at this level of detail.
\item \textbf{Strehl ratio and FWHM are both confirmed unstable/non-monotonic} in this
  project's near-field operating regime (Chapter~\ref{ch:mtf_strehl}); neither should
  be used as a primary design metric without further work.
\item \textbf{The $6.25$--$9\%$ brightness floor (Chapter~\ref{ch:stage3a}) was
  reverse-engineered to match the known clinical range, not independently derived.}
  A separate absolute-threshold analysis (Chapter~\ref{ch:brightness_reconciliation})
  exists but is not yet fully reconciled against it.
\item \textbf{No acuity predictor --- closed-form or ideal-observer, calibrated or
  uncalibrated --- has passed validation against real clinical data}
  (Chapter~\ref{ch:validation}). This is confirmed, not assumed, and is diagnosed as
  predominantly (though evidence suggests not entirely) a data-scarcity problem: more
  and better public clinical data, not further model tuning, is what this needs.
\item \textbf{The Watson--Yellott pupil--illuminance feedback coupling used throughout
  Stage 4} (Chapter~\ref{ch:stage4}) is a first-principles construction, not itself
  drawn from a single cited source; whether real pupils fully re-equilibrate to it
  within a typical clinical test's timescale is untested.
\item \textbf{This entire report is theory and numerical simulation only.} No physical
  prototype or human-subject experiment has been built or run against any prediction
  in this document.
\end{enumerate}

\section{Suggested priorities for closing these gaps}

Roughly in order of cost vs.\ payoff, per this project's own standing assessment:
(a) reconciling item 6 against item 3's absolute thresholds is the cheapest remaining
task --- both pieces of code already exist and simply need to be run against each
other at a shared, explicitly stated scene-luminance assumption; (b) designing and
running an actual physical/clinical experiment, with per-subject refractive error
recorded (the one fixable flaw every public dataset used in this report shares), would
directly address item 7, the single highest-value open item; (c) further acuity-model
refinement without new data is, per Chapter~\ref{ch:validation}'s own evidence, less
likely to move the outcome than either of the above.
